\chapter{Auswertung}
\label{ch:Auswertung}
% Liste der genutzer Formeln für die Fehlerrechnung
\section*{Fehlerrechnung}
Für die statistische Auswertung von $n$ Messwerten $x_i$ werden folgende Größen definiert \cite{errorSkript25}:
\begin{align}
    \bar{x} &= \frac{1}{n} \sum_{i=1}^{n} x_i \vphantom{\sqrt{\sum_i^n}^2} && \text{\textcolor{gray}{Arithmetisches Mittel}} \label{eq:arithmetisches_mittel} \\
    \sigma^2 &= \frac{1}{n-1} \sum_{i=1}^{n} (x_i - \bar{x})^2 \vphantom{\sqrt{\sum_i^n}^2} && \text{\textcolor{gray}{Variation}} \label{eq:variation} \\
    \sigma &= \sqrt{\frac{1}{n-1} \sum_{i=1}^{n} (x_i - \bar{x})^2} \vphantom{\sqrt{\sum_i^n}^2} && \text{\textcolor{gray}{Standardabweichung}} \label{eq:standardabweichung} \\
    \Delta \bar{x} &= \frac{\sigma}{\sqrt{n}} = \sqrt{\frac{1}{n(n-1)} \sum_{i=1}^n(\bar x - x_i)^2} \vphantom{\sqrt{\sum_i^n}^2} && \text{\textcolor{gray}{Fehler des Mittelwerts}} \label{eq:fehler_mittelwert} \\
    \Delta f &= \sqrt{\left(\frac{\partial f}{\partial x} \Delta x\right)^2 + \left(\frac{\partial f}{\partial y} \Delta y\right)^2} \vphantom{\sqrt{\sum_i^n}^2} && \text{\textcolor{gray}{Gauß’sches Fehlerfortpflanzungsgesetz für $f(x,y)$}} \label{eq:gauss_fehlfortpflanzung} \\
    \Delta f &= \sqrt{(\Delta x)^2 + (\Delta y)^2} \vphantom{\sqrt{\sum_i^n}^2} && \text{\textcolor{gray}{Fehler für $f = x + y$}} \label{eq:fehler_summe} \\
    \Delta f &= |a| \Delta x \vphantom{\sqrt{\sum_i^n}^2} && \text{\textcolor{gray}{Fehler für $f = ax$}} \label{eq:fehler_proportional} \\
    \frac{\Delta f}{|f|} &= \sqrt{\left(\frac{\Delta x}{x}\right)^2 + \left(\frac{\Delta y}{y}\right)^2} \vphantom{\sqrt{\sum_i^n}^2} && \text{\textcolor{gray}{relativer Fehler für $f = xy$ oder $f = x/y$}} \label{eq:relativer_fehler} \\
    \sigma &= \frac{|a_{lit} - a_{gem}|}{\sqrt{\Delta a_{lit}^2 + \Delta a_{gem}^2}} \vphantom{\sqrt{\sum_i^n}^2} && \text{\textcolor{gray}{Berechnung der signifikanten Abweichung}} \label{eq:signifikante_abweichung}
\end{align}

\newpage

\section{Visuallisierung der Teilchenbewegung}
\label{sec:A1}
Im ersten Aufgabenteil soll der zurückgelegte Weg des beobateten Tilchens geplottet werden, um einen besseren Eindruck der BEwegung zu bekommen.
Die Bewegung ist der \cabb{Bewegung-1a} zu entnehmen. Zwischen zwei Positionen liegt ungefähr eine Sekunde.

\img{Bewegung-1a}{Relative Bewegung des beobachteten Teilchens; Start und Ende vermerkt. Es wurden 150 Messpunkte aufgenommen.}

\section{Normalverteilung und Verbindung}
\label{sec:A2}
Weiterführend soll untersucht werden, ob unsere Messdaten stringent mit der Theorie sind. Insbesondere bedeutet dies, dass geprüft wird, ob die mittlere Verschiebung pro Zeitintervall eine Gauß-Normalverteilung darstellt.
Für jede Sekunde wurde die x- und y-Position gemessen. Aufgrund dessen, dass die brownsche Bewegung als isotrop angenommen wird, werden diese Werte zu einem Wert verrechnet.
Für die Gaußverteilung muss das \hyperref[eq:arithmetisches_mittel]{arithmetische Mittel (\ref*{eq:arithmetisches_mittel})} $\mu$ und die \hyperref[eq:standardabweichung]{Standardabweichung (\ref*{eq:standardabweichung})} $\sigma$ bestimmt werden.
Diese Theoretisch erwartete Verteilung wird über das Histogramm gelegt, um eine Abschätzung zu treffen, ob die theoretische Gaußkruve tatsächliche die Messwerte wiederspiegelt. 
In \cabb{Histogramm-2b} sind zwei Histogramme verschiedener Bin-Breiten eingestellt. 

\img[0.95]{Histogramm-2b}{Histogramm zur Untersuchung der Normalverteilung mit Gaußverteilung. Links mit 15 Bins, rechts mit 50. \afz{Breite} meint dabei das Gesamtintervall der Abszisse.}

Der rechte Plot mit 15 Bins zeigt eine ziemlich gute Nährung der Gauskruve an die Verteilung des Histogramms. Verkleinert man jedoch die Bin-Breite auf 50 Bins, so fällt ein Peak rund um $0$ auf. 
Dabei sind kleine Verschiebungen bei $\approx 0$ statistisch überpräsentiert. Hier soll nochmal in der Diskussion (\cch{Diskussion}) drauf eingegangen werden. 
Allgmeien lässt sich dennoch sagen, dass die gemessenen Werte der Theorie schon nah kommen, sodass die Annahme hält, dass es sich um die bronsche Bewegung handelt. 

\section{Bestimmung der Boltzmannkonstante}
\label{sec:A3}
Da die bronwsche Bewegung als wahr angenommen werden kann, kann weitergehend die Bolzmannkonstante bestimmt werden. 

\subsection{Bestimmung über das mittlere Verschiebungsquadrat}
\label{ssec:A3_1}
Zur Bestimmung der Boltsmannkonstante $k_b$ kann das mittlere Verschiebungsquadrat $\langle r^2 \rangle$. Dieses berechnet sich über die arithmetische Mittel der quadrierten Bewegungen in x- und y-Richtung:

\eq{r_quadr}{
    \langle r^2 \rangle = \frac{1}{149} \sum_{i=}^{149}({x_i}^2 + {y_i}^2) \, .
}

Die Ungenauigkeit ist dabei über den \hyperref[eq:fehler_mittelwert]{Fehler des Mittelwertes (\ref*{eq:fehler_mittelwert})} definiert,

\eq{err_r_quadr}{
    \Delta \langle r^2 \rangle = \frac{\sigma}{\sqrt{149}}
}

mit der Standardabweichung $\sigma$. b